\subsection{Transport affine des coefficients de Laurent}
\label{ix-add:stieltjes}

Le changement positif et adimensionnel $Q\mapsto\lambda Q$, avec
$\lambda>0$, agit sur le lecteur ultérieur au moyen de
\[
 Z_\lambda(s)=\lambda^{-s}Z(s).
\]
L'identité se vérifie d'abord dans $\Re s>1$ et se prolonge
méromorphiquement. Le résidu de $Z_\lambda$ en un est $\lambda^{-1}$.
Pour comparer les coefficients avec un résidu unitaire, on définit
\[
 \widehat Z_\lambda(s)=\lambda Z_\lambda(s)
 =e^{-\ell(s-1)}Z(s),\qquad \ell=\log\lambda,
\]
\[
 \widehat Z_\lambda(s)=\frac1{s-1}
 +\sum_{n\geq0}\frac{(-1)^n\gamma_n(\lambda)}{n!}(s-1)^n.
\]

\HMTresultado{Théorème}{Loi de transport de Stieltjes}{ix-add:stieltjes-law}
Pour tout $n\geq0$,
\[
 \gamma_n(\lambda)=
 \sum_{j=0}^{n}\binom nj\gamma_j\ell^{n-j}
 -\frac{\ell^{n+1}}{n+1}.
\]
Cette transformation est triangulaire affine. Deux changements d'échelle se composent en multipliant leurs paramètres et l'échelle inverse fournit la transformation inverse.

\textbf{Démonstration.} Posons $z=s-1$ et multiplions la série
complète de Laurent par $e^{-\ell z}$. Le terme polaire contribue au
coefficient de $z^n$ par $(-\ell)^{n+1}/(n+1)!$ ; les termes réguliers
contribuent par
\[
 \sum_{j=0}^n\frac{(-1)^j\gamma_j}{j!}
                    \frac{(-\ell)^{n-j}}{(n-j)!}.
\]
Multiplier la somme par $(-1)^nn!$ produit la formule. Le produit
$e^{-\ell z}e^{-mz}=e^{-(\ell+m)z}$ prouve la loi de composition.
La substitution $m=-\ell$ donne l'inverse. \hfill$\square$

En particulier,
\[
 \gamma_0(\lambda)=\gamma_0-\ell,\qquad
 \gamma_1(\lambda)=\gamma_1+\gamma_0\ell-\frac{\ell^2}{2}.
\]
Le terme affine enregistre la contribution du pôle et intervient dès le premier coefficient. Pour $\lambda=3$, on retrouve la transformation de la classe résiduelle zéro : ses coefficients sont ceux de $Z_0(s)=3^{-s}Z(s)=\widehat Z_3(s)/3$ de la section~\ref{lectura:manuscrito-04}. Cette spécialisation conserve le facteur $1/3$ du résidu. La loi générale permet de transporter la normalisation sans recalculer ni sélectionner un autre lecteur primaire TPK.
